\subsection{Extension of a linear representation of G to the measures on G}

Let G be a locally compact group, E a locally convex space, U a linear representation of G in E. Assume U to be continuous and E quasi-complete. Then, for every measure \(\mu \in \mathcal{C}'(G)\) , one has

\[
\int_ {\mathbf {G}} U (s) d \mu (s) \in \mathcal {L} (\mathbf {E}; \mathbf {E})
\]

(Ch. VI, §1, No.~7). We shall write \(U(\mu) = \int_{\mathbf{G}} U(s) d\mu(s)\) . We equip \(\mathcal{C}'(\mathbf{G})\) with the topology of compact convergence in \(\mathcal{C}(\mathbf{G})\) . The mapping \((\mu, x) \mapsto U(\mu)x\) of \(\mathcal{C}'(\mathbf{G}) \times \mathbf{E}\) into \(\mathbf{E}\) is hypocontinuous relative to the equicontinuous subsets of \(\mathcal{C}'(\mathbf{G})\) and the compact subsets of \(\mathbf{E}\) ; in particular, the mapping \(\mu \mapsto U(\mu)\) of \(\mathcal{C}'(\mathbf{G})\) into \(\mathcal{L}(\mathbf{E}; \mathbf{E})\) (equipped with the topology of compact convergence) is continuous (loc. cit., Prop. 16).

In order to be able to apply these results later on, we note that if X is a locally compact space then \(\mathcal{C}(\mathrm{X})\) , equipped with the topology of compact convergence, is complete (GT, X, §1, No.~6, Cor. 3 of Th. 2). On the other hand, \(\mathcal{K}(\mathrm{X})\) is barreled, therefore its dual \(\mathcal{M}(\mathrm{X})\) , equipped with the topology of compact convergence in \(\mathcal{K}(\mathrm{X})\) , is quasi-complete (TVS, III, §4, No.~2, Cor. 4 of Th. 1). Of course, \(\overline{\mathcal{K}(\mathrm{X})}\) is complete for the topology deduced from its norm, therefore its dual \(\mathcal{M}^{1}(\mathrm{X})\) is quasi-complete for the topology of compact convergence in \(\overline{\mathcal{K}(\mathrm{X})}\) (loc. cit.).

Let us now assume that U is a continuous linear representation of the locally compact group G on a Banach space E. Set \(g(s) = \|U(s)\|\) for all \(s \in G\) . Then, if \(\mu\) is a measure on G such that g is \(\mu\) -integrable, one has \(\int_{\mathbf{G}} U(s) d\mu(s) \in \mathcal{L}(\mathbf{E}; \mathbf{E})\) and \(\left\|\int_{\mathbf{G}} U(s) d\mu(s)\right\| \leqslant \int g(s) d|\mu|(s)\) (Ch. VI, §1, No.~7, Remark 1). We again write \(U(\mu) = \int_{\mathbf{G}} U(s) d\mu(s)\) .

\begin{enumerate}
\def\labelenumi{\arabic{enumi}.}
\setcounter{enumi}{6}
\tightlist
\item
  Relations between the endomorphisms \(U(\mu)\) and the endomorphisms \(U(s)\)
\end{enumerate}

Lemma 4. --- Let T be a locally compact space, a a point of T, M a subset of \(\mathcal{M}(\mathrm{T})\) , and \(\mathfrak{F}\) a filter on M. Assume that:

\begin{enumerate}
\def\labelenumi{(\roman{enumi})}
\item
  for every compact subset \(\mathbf{K}\) of \(\mathbf{T}\) , the numbers \(|\mu|(\mathbf{K})\) , for \(\mu \in \mathbf{M}\) , are bounded above;
\item
  \(\lim_{\mu, \mathfrak{F}} |\mu| (\mathrm{K}) = 0\) for every compact subset \(\mathrm{K}\) of \(\mathrm{T} - \{a\}\) .
\item
  there exists a compact neighborhood V of a in T such that \(\lim_{\mu, \mathfrak{F}} \mu(V) = 1\) .
\end{enumerate}

Then the filter \(\mathfrak{F}\) converges to \(\varepsilon_{a}\) in \(\mathcal{M}(\mathrm{T})\) equipped with the topology of compact convergence in \(\mathcal{H}(\mathrm{T})\) .

By the hypothesis (i), M is an equicontinuous subset of \(\mathcal{M}(\mathrm{T})\) since it is vaguely bounded and \(\mathcal{H}(\mathrm{T})\) is barreled (TVS, III, §4, No.~2, Th. 1). It therefore suffices (GT, X, §2, No.~4, Th. 1) to prove that if \(f \in \mathcal{H}(\mathrm{T})\) , then \(\lim_{\mu, \mathfrak{F}} \mu(f) = f(a)\) . Let K be the union of V and the support of \(f\) ; if \(\mathrm{K}'\) is the closure of \(\mathrm{K} - \mathrm{V}\) , one has

\[
| \mu (\mathrm{K}) - \mu (\mathrm{V}) | = | \mu (\mathrm{K} - \mathrm{V}) | \leqslant | \mu | (\mathrm{K} ^ {\prime});
\]

since \(\mathbf{K}'\) is compact and does not contain \(a\) , one concludes from this that \(\lim_{\mu, \mathfrak{F}} \mu(\mathbf{K}) = 1\) . Let \(\varepsilon > 0\) , and let \(\mathbf{W}\) be an open neighborhood of \(a\) in \(\mathbf{K}\) such that \(|f(t) - f(a)| \leqslant \varepsilon\) for \(t \in \mathbf{W}\) ; one can write

\[
\mu (f) - f (a) = f (a) \big (\mu (\mathrm{K}) - 1 \big) + \int_ {\mathrm{K}} \big (f (t) - f (a) \big) d \mu (t);
\]

the integral over K may be written as the sum of the analogous integrals over W and K - W; if \(C = \sup |f|\) , one therefore has

\[
| \mu (f) - f (a) | \leqslant \mathrm{C} | \mu (\mathrm{K}) - 1 | + \varepsilon \cdot | \mu | (\mathrm{K}) + 2 \mathrm{C} \cdot | \mu | (\mathrm{K} - \mathrm{W}).
\]

Since the first and third terms on the right side tend to 0 with respect to \(\mathfrak{F}\) , one sees that indeed \(\lim_{\mu \to \mathfrak{F}} \mu(f) = f(a)\) .

COROLLARY 1. --- With hypotheses as in Lemma 4, suppose in addition that there exists a compact subset \(K_{0}\) of T containing the supports of all of the measures \(\mu \in M\) . Then F also converges to \(\varepsilon_{a}\) in \(\mathcal{C}'(T)\) equipped with the topology of compact convergence in \(\mathcal{C}(T)\) .

For, the restriction mapping of \(\mathcal{C}(\mathrm{T})\) into \(\mathcal{C}(\mathrm{K}_{0})\) is continuous; therefore, if H is a compact subset of \(\mathcal{C}(\mathrm{T})\) , then the restrictions to \(K_{0}\) of the functions in H form a compact subset of \(\mathcal{C}(\mathrm{K}_{0})\) . It then suffices to apply Lemma 4 on replacing T by \(K_{0}\) .

COROLLARY 2. --- With hypotheses as in Cor. 1, let f be a continuous mapping of T into a quasi-complete locally convex space E. Then

\[
\lim _ {\mu , \mathfrak {F}} \int f (t) d \mu (t) = f (a).
\]

This follows from Cor. 1, and Prop. 14 of Ch. VI, §1, No.~6.

COROLLARY 3. --- Let G be a locally compact group, E a quasi-complete locally convex space, and U a continuous linear representation of G in E. Let \(\beta\) be a positive measure on G, a an element of G, and B a base for the filter of neighborhoods of a, formed of compact neighborhoods. For every \(V \in B\) , let \(f_{V}\) be a continuous function \(\geqslant 0\) on G, with support contained in V, and such that \(\int f_{V} d\beta = 1\) . Then, for every \(x \in E\) ,

\[
U (a) x = \lim _ {\mathrm{V}} U (f _ {\mathrm{V}} \cdot \beta) x,
\]

the limit being taken with respect to the section filter of \(\mathfrak{B}\) .

The mapping \(s \mapsto U(s)x\) of \(G\) into \(E\) is continuous. By Cor. 2, \(U(a)x = \lim_{V} \int (U(s)x) \cdot f_V(s) d\beta(s)\) with respect to the section filter of \(\mathfrak{B}\) , that is, \(U(a)x = \lim_{V} U(f_V \cdot \beta)x\) .

PROPOSITION 10. --- Let G be a locally compact group, E a quasi-complete locally convex space, U a continuous linear representation of G in E, and \(\beta\) a positive measure on G with support G.

\begin{enumerate}
\def\labelenumi{(\roman{enumi})}
\item
  The vectors \(U(f \cdot \beta)x\) , where \(f\) runs over \(\mathcal{K}(\mathbf{G})\) and \(x\) runs over \(\mathbf{E}\) , are dense in \(\mathbf{E}\) .
\item
  Let \(\mathbf{F}\) be a closed linear subspace of \(\mathbf{E}\) . If \(\mathbf{F}\) is stable for \(U\) , then \(U(\mu)(\mathbf{F}) \subset \mathbf{F}\) for every \(\mu \in \mathcal{C}'(\mathbf{G})\) . Conversely, if \(U(f \cdot \beta) \subset \mathbf{F}\) for every \(f \in \mathcal{K}(\mathbf{G})\) , then \(\mathbf{F}\) is stable for \(U\) .
\end{enumerate}

The first part of (ii) is immediate, since the restrictions of the \(U(s)\) to F ( \(s \in G\) ) define a continuous linear representation of G in the quasi-complete locally convex space F. The second part of (ii), and (i), follow from Cor. 3 of Lemma 4.

